\subsection{Approximate complete intersections}
\label{aci:section}

We prove Walsh's approximate complete-intersection theorem
\cite[Theorem~1.3]{walsh2020}. The partial-degree framework and the
construction of compatible cuts follow
\cite[Definition~2.2 and Lemma~5.3]{walsh2020}.
The comparison of restriction spaces appears in the proof of
\cite[Lemma~3.5]{walsh2020}. Walsh obtains a corresponding prefix
degree estimate for his chosen cuts
\cite[Theorem~5.5]{walsh2020}.
We prove the forms needed here; the component Hilbert bounds in
Lemma~\ref{poly:hilbert} are the only quantitative algebraic input.
In particular, this proof is independent
of Theorem~\ref{conc:theorem}.

\begin{theorem}[Approximate complete intersection]
\label{aci:theorem}
\label{conc:walsh}
For every $n$ there is a constant $C_n$ with the following property.
If $V\subset\mathbb A^n_{\mathbb C}$ is irreducible of codimension
$c>0$, there are nonconstant polynomials $f_1,\ldots,f_c\in I(V)$ such
that $V$ is an irreducible component of $Z(f_1,\ldots,f_c)$ and
\[
 \prod_{i=1}^c\deg f_i\le C_n\deg V.
\]
\end{theorem}

Here and below, a component \emph{containing $V$} means that it contains
the whole irreducible variety $V$.  Set
\[
 I_{\le t}(V)=I(V)\cap\mathbb C[x_1,\ldots,x_n]_{\le t},
 \qquad
 H_V(t)=\dim_{\mathbb C}
 \bigl(\mathbb C[x_1,\ldots,x_n]_{\le t}/I_{\le t}(V)\bigr).
\]
For $1\le i\le c$, define $\delta_i$ to be the smallest integer $t$
such that every irreducible component of $Z(I_{\le t}(V))$ containing
$V$ has codimension at least $i$.  The notation $Z(I_{\le t}(V))$
means the zero set of a basis of this finite-dimensional vector space.
Finite generation of $I(V)$ shows that these integers exist, and their
definition gives
\begin{equation}
 1\le\delta_1\le\cdots\le\delta_c.
 \label{aci:partial-degrees}
\end{equation}
We also put $\delta_0=0$.

\begin{lemma}[Compatible partial cuts]
\label{aci:cuts}
There are $f_i\in I_{\le\delta_i}(V)$, $1\le i\le c$, such that
every component of $Z(f_1,\ldots,f_i)$ containing $V$ has codimension
exactly $i$.

For any such choice, if $0\le j<c$ and
$\delta_j\le m<\delta_{j+1}$, some component $W$ of
$Z(f_1,\ldots,f_j)$ containing $V$ satisfies
\begin{equation}
 I_{\le m}(W)=I_{\le m}(V),\qquad H_W(m)=H_V(m).
 \label{aci:shared-rank}
\end{equation}
When $j=0$, the prefix zero set and its unique component are
$\mathbb A^n$.
\end{lemma}

\begin{proof}
Suppose $f_1,\ldots,f_{i-1}$ have been chosen, and list the finitely
many components $W$ of their zero set which contain $V$.  Each has
codimension $i-1$.  None is contained in $Z(I_{\le\delta_i}(V))$:
otherwise a component of that latter set would contain $W$, hence
contain $V$ and have codimension at most $i-1$, contrary to the
definition of $\delta_i$.  Consequently the members of
$I_{\le\delta_i}(V)$ that vanish on any fixed $W$ form a proper
linear subspace.  Choose $f_i$ outside their finite union.  It cuts
every such $W$ properly.  The principal ideal theorem then gives
codimension exactly $i$ for every new component containing $V$.
This includes $i=1$, when the preceding zero set is the ambient space,
and proves the construction.

For the second assertion, suppose no component $W$ in the indicated
prefix has $I_{\le m}(V)\subset I(W)$.  The same finite-union argument
provides $g\in I_{\le m}(V)$ which vanishes identically on none of
them.  All components of $Z(f_1,\ldots,f_j,g)$ containing $V$ then
have codimension $j+1$.  Since $m\ge\delta_j$, every one of these
equations belongs to $I_{\le m}(V)$.  Thus every component of
$Z(I_{\le m}(V))$ containing $V$ has codimension at least $j+1$.
This contradicts $m<\delta_{j+1}$.  For a component $W$ where the
inclusion does hold, the reverse inclusion follows from $V\subset W$,
giving \eqref{aci:shared-rank}.
\end{proof}

\begin{proof}[Proof of Theorem~\ref{aci:theorem}]
Fix the polynomials from Lemma~\ref{aci:cuts}.  We will prove, for
$0\le s\le c$, that every component $U$ of the first $s$ cuts
containing $V$ satisfies
\begin{equation}
 \deg U\ge\gamma_n^{\,s}\prod_{i=1}^s\delta_i
 \label{aci:prefix-degree}
\end{equation}
for one sufficiently small $0<\gamma_n\le1$ depending only on $n$.
For $s=0$ both sides are $1$.

We record the precise consequence of Lemma~\ref{poly:hilbert} used
here.  There are constants $a_n>0$ and $b_n\ge1$ such that an
irreducible $k$-dimensional variety $X$ satisfies
\begin{equation}
 H_X(m)\le b_n\deg(X)m^k\qquad(m\ge1),
 \label{aci:rank-upper}
\end{equation}
and, if $X$ is a component of a system whose equations have degrees at
most $m$, then
\begin{equation}
 H_X(m)\ge a_n\deg(X)m^k\qquad(m\ge1).
 \label{aci:rank-lower}
\end{equation}
The second assertion is the bound at the equation degree in
\eqref{poly:hilbert-at-L}, applied with $L=m$.  It permits other
components of the system to have larger dimension.  The lower bound
also holds for $X=\mathbb A^n$, after decreasing $a_n$, directly from
$H_{\mathbb A^n}(m)=\binom{m+n}{n}$.

Suppose \eqref{aci:prefix-degree} is known for smaller indices, and
let $U$ be a component at index $s$.  If $\delta_s=1$, the entire
product in \eqref{aci:prefix-degree} is $1$, and $\deg U\ge1$ proves
the claim.  Otherwise set
\[
 m=\delta_s-1,\qquad
 j=\max\{i:0\le i<s,\ \delta_i<\delta_s\}.
\]
Thus
\begin{equation}
 \delta_j\le m<\delta_{j+1}
   =\cdots=\delta_s\le2m.
 \label{aci:block}
\end{equation}
Lemma~\ref{aci:cuts} supplies a component $W$ at index $j$ with
$H_W(m)=H_V(m)$.  Its dimension is $n-j$, and the degrees of its
defining prefix equations are at most $\delta_j\le m$.
The induction hypothesis and \eqref{aci:rank-lower} therefore give
\[
 H_V(m)=H_W(m)
 \ge a_n\gamma_n^{\,j}
       \left(\prod_{i=1}^j\delta_i\right)m^{n-j}.
\]
This formula also applies when $j=0$.  Since $V\subset U$, restriction
from $U$ to $V$ is surjective on polynomial restriction spaces, so
$H_U(m)\ge H_V(m)$.  On the other hand, $\dim U=n-s$ and
\eqref{aci:rank-upper} bounds $H_U(m)$ above.  Combining the bounds,
then using \eqref{aci:block}, yields
\[
 \begin{split}
 \deg U
 &\ge \frac{a_n}{b_n}\gamma_n^{\,j}
       \left(\prod_{i=1}^j\delta_i\right)m^{s-j}\\
 &\ge \frac{a_n}{2^n b_n}\gamma_n^{\,j}
       \prod_{i=1}^s\delta_i.
 \end{split}
\]
Choose $\gamma_n\le\min\{1,a_n/(2^n b_n)\}$.  Since $j<s$, the
last expression is at least the right side of
\eqref{aci:prefix-degree}, completing the induction.  Notice that no
containment relation between $U$ and the chosen $W$ was needed; both
contain $V$.

At $s=c$, every prefix component containing $V$ has the same dimension
as $V$, so $V$ itself is such a component.  Applying
\eqref{aci:prefix-degree} to it gives
\[
 \prod_{i=1}^c\deg f_i
 \le\prod_{i=1}^c\delta_i
 \le\gamma_n^{-c}\deg V
 \le\gamma_n^{-n}\deg V.
\]
This proves the theorem with $C_n=\gamma_n^{-n}$.
\end{proof}

The conclusion concerns an irreducible component.  It does not assert
that the displayed equations generate its radical ideal or have
independent differentials; the next regularity argument supplies the
additional control needed here.
